\ifdefined\MyWebBook\else
\documentclass[../textbook.tex]{subfiles}
\begin{document}
\fi
\bookchapter{批处理调度}{ch:batching-scheduling}

并发请求的完成时间由算子执行与资源调度共同决定。服务系统还必须决定何时接纳请求、哪些请求共同执行、长提示如何穿插，以及资源不足时谁应等待。调度改变的是请求共享资源的方式，其收益与到达过程、长度分布和服务目标有关。

\bookxref{ch:incremental-decoding-cache}给出可复用状态及其容量，\bookxref{ch:inference-algorithms}给出一次执行的成本。本章将二者放入随时间变化的工作负载中。部署、连接协议与实例生命周期见\bookxref{ch:serving-architecture}。

\section{服务性能的计量约定}

本章主要符号见表\ref{tab:batching-scheduling-symbols}。
\begin{table}[htbp]
\centering
\caption{批处理调度主要符号}\label{tab:batching-scheduling-symbols}
\begin{tabular}{ll}
\toprule
符号 & 含义\\
\midrule
$a_i,b_i,c_i$ & 请求 $i$ 到达、首次获得执行和终止的时间。\\
$u_{i,j}$ & 客户端收到请求 $i$ 的第 $j$ 个输出词元的时间。\\
$P_i,O_i$ & 输入长度与实际输出长度；$O_i^{\max}$ 为接纳时约定的输出上限。\\
$W_i,T_i,S_i$ & 首次执行前等待时间、系统逗留时间、理想队列模型中的服务时间。\\
$N(t),N_q(t)$ & 时刻 $t$ 系统内和等待队列内的请求数。\\
$\lambda,\mu,\rho$ & 队列模型中的到达率、单服务台服务率和利用率。\\
$B,K,M$ & 单轮新计算词元预算、活动序列上限和可用缓存字节预算。\\
$q_i,s_i,m_i$ & 单轮为请求分配的新位置数、已有上下文长度及缓存字节数。\\
$d_i,w_j$ & 请求截止时间与租户 $j$ 的调度权重。\\
\bottomrule
\end{tabular}
\end{table}

\term{首词元时延}{Time to First Token}{TTFT}与\term{词元间时延}{Inter-Token Latency}{ITL}分别为
\begin{equation}
 \begin{aligned}
  \operatorname{TTFT}_i&=u_{i,1}-a_i,\\
  \operatorname{ITL}_{i,j}&=u_{i,j}-u_{i,j-1}\quad(j\ge2).
 \end{aligned}
\end{equation}
请求生成至少两个词元时，其\term{平均输出词元时间}{Time per Output Token}{TPOT}定义为
\begin{equation}
 \operatorname{TPOT}_i=\frac{u_{i,O_i}-u_{i,1}}{O_i-1}.
 \label{eq:schedule-tpot}
\end{equation}
因此最后一个词元到达的时延恰为 $\operatorname{TTFT}_i+(O_i-1)\operatorname{TPOT}_i$。这里的 TTFT 已包含到首次词元之间的排队，再叠加 $W_i=b_i-a_i$ 会重复计入。只有一个词元时 TPOT 不定义；无输出的失败请求计入失败统计，其输出时延记为未定义。

客户端时间还包含网络和缓冲，工作进程内的执行时间不包含这些部分。二者都可记录，但必须说明观测位置。流式协议可能一次传输多个词元；若只能观察数据块，相应指标应称为块间隔，它与逐词元 ITL 不同。\footnote{同一请求的终止控制消息可能晚于最后一个词元；因此完整协议时延 $c_i-a_i$ 还可包括收尾开销。}

吞吐也有不同口径：请求完成数、输入词元数、输出词元数和满足服务目标的完成数。预填充与解码的计算和搬运成本不同，$\lambda\mathbb E[P_i+O_i]$ 仅表示词元流量，与实际设备需求之间隔着一层换算。缓存命中又会使输入长度与实际计算长度不同。

\section{连续批处理}

\subsection{静态批处理的容量浪费}
\term{静态批处理}{Static Batching}{}在批次初始化时锁定参与请求集合，通常直至当前批次全部序列生成结束方才接纳下一组请求。在静态批处理中，请求尽管可以在遇到终止符时立即向客户端流式返回，但在整批全部请求结束前，已被释放的显存槽位与计算单元无法被新到达的请求动态填补，这种生命周期木桶效应直接导致硬件利用率随生成步数推进而显著衰减。

\begin{example}
设调度模型如下：所有预填充均已完成，一轮解码固定耗时 \qty{10}{\milli\second}，最多容纳两条并发序列。三个请求同时到达，剩余输出长度分别为 $A=4,B=1,C=1$。静态批处理先组批执行 $A,B$：$B$ 在 \qty{10}{\milli\second} 完成并退出，$A$ 在 \qty{40}{\milli\second} 完成，随后调入执行 $C$ 并在 \qty{50}{\milli\second} 完成；三个请求的平均完成时延为 $\frac{10+40+50}{3}\approx \qty{33.33}{\milli\second}$。若采用动态调度，在第2轮即时将 $C$ 填补入 $B$ 释放的空闲槽位，完成时延分别优化为 \qty{40}{\milli\second}、$B$ 为 \qty{10}{\milli\second}、$C$ 为 \qty{20}{\milli\second}，平均完成时延降至 $\frac{40+10+20}{3}\approx \qty{23.33}{\milli\second}$。

\begin{figure}[htbp]
\centering
\begin{tikzpicture}[x=1.15cm,y=.65cm,font=\small]
\node[anchor=east] at (0,2) {静态：槽1};
\node[anchor=east] at (0,1) {静态：槽2};
\node[anchor=east] at (0,-1) {连续：槽1};
\node[anchor=east] at (0,-2) {连续：槽2};
\foreach \x in {0,...,3}{
 \draw[fill=blue!15] (\x,1.65) rectangle +(1,.7);
 \node at (\x+.5,2) {A};
 \draw[fill=blue!15] (\x,-1.35) rectangle +(1,.7);
 \node at (\x+.5,-1) {A};
}
\draw[fill=green!15] (0,.65) rectangle +(1,.7);
\node at (.5,1) {B};
\draw[fill=orange!20] (4,1.65) rectangle +(1,.7);
\node at (4.5,2) {C};
\draw[fill=green!15] (0,-2.35) rectangle +(1,.7);
\node at (.5,-2) {B};
\draw[fill=orange!20] (1,-2.35) rectangle +(1,.7);
\node at (1.5,-2) {C};
\draw[->] (0,-3) -- (5.4,-3) node[right] {轮次};
\foreach \x in {1,...,5}{\node at (\x-.5,-2.75) {\x};}
\end{tikzpicture}
\caption[连续批处理动态迭代级调度]{连续批处理动态迭代级调度。在迭代步进调度下，新到达请求可及时复用已完成请求释放的槽位，消除静态批处理的等待空泡。}\label{fig:schedule-continuous}
\end{figure}

\term{连续批处理}{Continuous Batching}{}在迭代边界调整活动成员。Orca 将这种决策落实为迭代级调度，并根据算子特性选择批处理方式\cite{yu2022orca}。真实系统的轮时会随批大小、上下文长度和算子形状变化，故图\ref{fig:schedule-continuous}只证明该构造中的收益，与吞吐预测无关。
\end{example}

表\ref{tab:batching-paradigms}比较了请求级批处理、逐步调整批次、分块预填充和抢占等调度方法。它们以不同粒度分配执行时间和缓存空间，延迟、吞吐与 SLO 达成率还取决于请求长度分布、到达率和设备容量。

\begin{table}[htbp]
\centering
\caption{推理批处理调度粒度与机制对比}\label{tab:batching-paradigms}
\begin{tabularx}{\linewidth}{@{}lXXXX@{}}
\toprule
调度范式 & 调度决策粒度 & 核心机制与状态管理 & 资源利用与延迟表现 & 典型适用场景与限制 \\
\midrule
静态批处理 & 批次完成级 & 请求填满固定批量后统一发射，全批完成方退出 & 填充浪费严重，利用率低；短请求受制于长请求 & 离线批量评估；在线高并发下极易拥塞 \\
动态聚合批处理 & 窗口超时级 & 在最大等待时间内聚批，批内各请求同步执行 & 相比静态减少排队，但仍存在尾部空转气泡 & 早期的专用服务引擎（如 Triton 传统后端） \\
连续批处理 & 单步迭代级 & 每次自回归迭代动态移出结束请求、插入就绪请求 & 消除尾部等待气泡，GPU 计算利用率提升数倍 & 现代通用 LLM 推理引擎（vLLM、SGLang、TGI） \\
分块预填充 & 词元预算级 & 将长 Prompt 拆分为等长小块，与解码请求共批执行 & 显著平抑长提示对解码 TPOT 的突发冲击 & 混合长上下文问答与高交互对话服务 \\
抢占与换出 & 显存告警级 & 显存耗尽时换出（Swap）或丢弃重算低优先级请求 & 避免 OOM 崩溃，严格保障高优先级 SLO & 突发高并发与长序列混合极端负载 \\
\bottomrule
\end{tabularx}
\end{table}

\subsection{静态批处理}
若内核把同批 $n$ 条序列填充到最大长度 $P_{\max}$，线性层与显式完整平方注意力矩阵的长度利用率分别为
\begin{equation}
 \begin{aligned}
  \eta_{\mathrm{linear}}&=\frac{\sum_iP_i}{nP_{\max}},\\
  \eta_{\mathrm{attention}}&=\frac{\sum_iP_i^2}{nP_{\max}^2}.
 \end{aligned}
\end{equation}
变长分块内核可跳过部分填充计算，实际利用率应按执行位置统计。

\term{长度分桶}{Length Bucketing}{}把相近长度放入同一执行组，减少形状差异。但等待相近请求会增加时延。动态聚批应同时设置最大等待时间与最大批量：任一条件触发就执行。仅为填满批次持续等待，会在低流量下恶化 TTFT。桶边界附近的长度抖动还可能造成低占用桶，应与\bookxref{ch:inference-algorithms}[中的形状分桶]中的编译缓存共同设计。

\section{分块预填充}

\subsection{分块预填充的上下文边界}
\term{分块预填充}{Chunked Prefill}{}将一个提示的新增位置分多轮处理，每块仍读取合法历史状态；它改变执行时序，不必改变模型可见上下文。SARATHI 研究了预填充分块与解码共同组批的执行方式\cite{agrawal2023sarathi}。具体分块长度须服从设备、批次和服务目标，论文中的某个配置只是对应那些条件的一组取值。

在已有 $s_i$ 个位置时处理 $q_i$ 个新位置，因果注意力所涉及的键查询配对数为
\begin{equation}
 q_i s_i+\frac{q_i(q_i+1)}2.
 \label{eq:schedule-pairs}
\end{equation}
推导是将第 $r$ 个新查询可见的 $s_i+r$ 个键对 $r=1,\ldots,q_i$ 求和。每块重读历史会产生搬运与启动开销；块越小，交互机会越多，总时间是否更短还取决于这些开销。解码通常每序列新增一个位置，预填充则可新增许多位置，因此一个简单预算是
\begin{equation}\label{eq:schedule-budgets}
 \begin{aligned}
  \sum_i q_i &\le B,\\
  |\mathcal A| &\le K,\\
  \sum_{p\in\mathcal P_{\rm live}}\operatorname{bytes}(p)+M_{\rm workspace} &\le M.
 \end{aligned}
\end{equation}
其中 $\mathcal P_{\rm live}$ 是实际驻留物理页集合，共享前缀只计一次。临时工作区要在可用容量中扣除，设备总显存与 KV 预算之间还差着权重和工作区。

仅限制 $B$ 仍不够：同样的新词元数在不同 $s_i$ 下有不同注意力成本。调度器还需使用经测量校准的批次成本估计 $\widehat t(\{q_i,s_i\})$，在交互目标允许的轮时内选择组合。该函数一般不同于各请求独立成本之和，因为共同执行会改变权重重用和有效带宽。

\begin{example}
\begin{assumption}
以下资源算例假设普通非共享 KV 缓存，每词元键值状态占用 \qty{64}{\kibi\byte}，单个物理页容纳 $16$ 个词元（即每物理页容量为 \qty{1}{\mebi\byte}）；设备显存扣除模型权重与运行时临时工作区后，分配给 KV 缓存的可用预算为 \qty{64}{\mebi\byte}（即 $64$ 个物理页）。实际缓存公式应依据网络层数、键值头数、数值精度及张量并行切分度确定。
\end{assumption}

系统当前有两条活动序列，历史长度分别为 $257$ 与 $255$ 个词元，分别占用 $\lceil 257/16\rceil=17$ 页与 $\lceil 255/16\rceil=16$ 页，合计占用 $33$ 页。若本轮为二者各解码一步，序列长度增长为 $258$ 与 $256$，物理页数依然维持 $33$ 页不变。此时调度器尝试并发处理一个新到达请求的 $256$ 词元前缀提示块，需要即时分配 $16$ 页，单步显存需求为 $33+16=49$ 页（满足 $\le 64$ 页预算）。然而，若该新请求声明的最大生成长度上限为 $1024$ 词元，其生命周期峰值需占用 $64$ 页，与既有序列峰值累加达 $97$ 页，远超 $64$ 页物理上限，三者无法全部无阻塞驻留至完成。

单轮执行仅需核验瞬时显存水位，而请求级准入控制则必须预估序列全生命周期的显存增长。调度系统既可采取悲观保守策略按最大长度上限预留显存，也可采取乐观超售策略允许动态申请，并在物理显存耗尽时依赖明确的挂起换出（Swap-out）或重算机制进行保底。基于分页映射的虚拟内存管理可从根本上消除外部内存碎片，并支持多请求间的公共系统前缀物理共享\cite{kwon2023pagedattention}。
\end{example}

\subsection{预填充与解码的优先级}
若策略始终绝对优先满足解码请求，可有效压低已接纳请求的词元间时延（ITL），但会导致排队中的新请求长时间得不到首次响应（TTFT 严重劣化）。反之，若始终抢占执行新到提示词元，正处于自回归阶段的长输出将被频繁挂起中断。\keyconcept{调度预填充偏向优化首词元时延（TTFT），调度增量解码偏向优化词元间时延（ITL）；分块预填充通过限定单轮新词元计算预算，实现两类异质负载的平滑时钟复用。}工业级引擎通常为预填充设定每轮计算词元预算（Chunk Budget）上限，在确保单轮解码迭代时延不超标的前提下分配剩余算力；多目标冲突时的调度策略权衡应通过细粒度的分位数指标加以度量。

\section{排队模型}

\subsection{Little 定律的推导}
Little 定律建立了系统稳态平均队长、请求到达率与平均逗留时间之间的恒等关系\cite{little1961queue}。设观察时间窗口为 $[0,H]$，系统内瞬时请求数表示为时间指示函数的求和：
\begin{equation}
 N(t)=\sum_i\mathbf1\{a_i\le t<c_i\}.
\end{equation}
对两端积分并交换求和次序：
\begin{equation}
 \int_0^H N(t)\,dt
 =\sum_i |[a_i,c_i)\cap[0,H]|.
 \label{eq:schedule-area}
\end{equation}
等式右侧为各请求在观察窗内的逗留时间累加。当系统处于统计稳态且长期平均到达率存在时，边界截断误差随 $H\to\infty$ 趋近于零。两端同除以 $H$ 并取极限，可得经典关系式：
\begin{equation}\label{eq:schedule-little}
 \begin{aligned}
  \overline N&=\lambda\mathbb E[T],\\
  \overline N_q&=\lambda\mathbb E[W_q].
 \end{aligned}
\end{equation}
其中第二式对等待队列区间重复应用相同面积分解所得。在连续批处理中，请求在迭代间歇可能多次处于等待显存的状态，$W_q$ 涵盖全周期的排队总耗时，大于首词元发射前的初始等待时间 $W_i$。

式\eqref{eq:schedule-little}对到达间隔与服务时间的分布形式无先验假设，但要求队长、吞吐率与时延基于同一样本统计边界。若计算到达率时计入了被准入控制拒绝的请求，却仅对成功接纳的请求统计驻留时间，等式将发生统计失真。例如某稳定服务节点平均每秒处理完成 $20$ 个请求，端到端平均驻留时间为 \qty{0.4}{\second}，则系统内平均在途请求数恒为 $\overline{N}=20\times 0.4=8$ 条。

\subsection{M/M/1 队列与长尾分布}
\begin{assumption}
假设请求到达服从强度为 $\lambda$ 的独立泊松过程，单处理单元的服务时间服从参数为 $\mu$ 的负指数分布，采用先到先服务（FCFS）且等待队列容量无上限，系统稳态服务利用率满足 $\rho=\frac{\lambda}{\mu}<1$。该连续时间马尔可夫链模型用于揭示排队延迟的非线性增长规律。
\end{assumption}

记系统内包含 $n$ 个请求的稳态概率为 $\pi_n$。由细致平衡条件（Detailed Balance）$\lambda\pi_n=\mu\pi_{n+1}$，解得几何分布形式 $\pi_n=(1-\rho)\rho^n$。系统内平均请求数与平均排队等待时间为
\begin{equation}\label{eq:schedule-mm1-mean}
 \begin{aligned}
  \mathbb E[N]&=\frac{\rho}{1-\rho},\\
  \mathbb E[T]&=\frac{1}{\mu-\lambda},\\
  \mathbb E[W_q]&=\frac{\rho}{\mu-\lambda}.
 \end{aligned}
\end{equation}
当利用率 $\rho\to 1$ 时，平均等待时延呈高阶双曲发散；高利用率以长尾排队时延的急速恶化为代价。

在 PASTA（Poisson Arrivals See Time Averages）性质下，新到达请求看到的系统状态精确等同于时间平均稳态分布。到达时系统中已有 $n$ 个在途请求，该请求完成需经历 $n+1$ 个参数为 $\mu$ 的独立指数服务阶段（首个请求的剩余服务时间由无记忆性保证）。对拉普拉斯变换求逆可得逗留时间的概率密度函数与分位数解析解：
\begin{equation}\label{eq:schedule-mm1-tail}
 \begin{aligned}
  \Pr(T>t)&=\conste{}^{-(\mu-\lambda)t},\\
  t_{0.99}&=\frac{\log 100}{\mu-\lambda}\approx \frac{4.605}{\mu-\lambda}.
 \end{aligned}
\end{equation}
若系统服务能力为 $\mu=100\ \text{req/s}$，在到达率 $\lambda=80\ \text{req/s}$ 下（$\rho=0.8$），平均时延为 \qty{50}{\milli\second}，P99 尾部时延约为 \qty{230}{\milli\second}；当到达率上升至 $\lambda=95\ \text{req/s}$ 时（$\rho=0.95$），平均时延激增至 \qty{200}{\milli\second}，P99 尾部时延恶化至 \qty{921}{\milli\second}。

\begin{note}[服务曲线的动态演化与测量基准]
在 LLM 推理场景中，多请求共享前向批量导致服务时间随在途批量、上下文长度及自回归步数非线性变化。理论排队模型中的服务率应以基准负载下的实测完成率为准。生产环境容量规划须采用基于实测轨迹的离散事件仿真（DES），结合实际输入输出长度分布校准长尾服务曲线。
\end{note}

\section{请求准入控制}

\term{准入控制}{Admission Control}{}发生在承诺接受任务的边界。至少要核查模型与租户配额、请求体与长度上限、缓存增长、队列空间和期限。可以用
\begin{equation}
 \widehat W_i+\widehat S_i+\Delta_i\le d_i-a_i
 \label{eq:schedule-admission}
\end{equation}
作为风险判断，其中 $\Delta_i$ 为估计误差和协议开销余量。分位数相加得不到同分位数的总和；要么直接估计联合总时延，要么使用明确的概率上界。\footnote{若 $\Pr(W>w)\le\epsilon_1$、$\Pr(S>s)\le\epsilon_2$，则不要求独立也有 $\Pr(W+S>w+s)\le\epsilon_1+\epsilon_2$。这来自并集界，与分位数可加性无关。}

\term{背压}{Backpressure}{}把下游容量不足传递给上游。它必须贯穿入口、有界等待队列、活动序列和输出缓冲。如果只限制设备并发，慢客户端仍可能使主机缓冲无限增长。若输出缓冲以每秒 $r_g$ 字节生成、每秒 $r_c$ 字节消耗，则积压近似按 $r_g-r_c$ 增长；当 $r_g>r_c$，容量 $H_b$ 的空缓冲将在 $\frac{H_b}{r_g-r_c}$ 秒后填满。持续消费不足时应暂停、取消或终止，并把资源释放传播到上游。

\begin{figure}[htbp]
\centering
\begin{tikzpicture}[node distance=5mm,font=\small,
 box/.style={draw,rounded corners,align=center,text width=23mm,minimum height=12mm}]
\node[box] (a) {入口\\长度与配额};
\node[box,right=of a] (b) {有界队列\\期限与公平};
\node[box,right=of b] (c) {逐轮执行\\词元与KV};
\node[box,right=of c] (d) {输出缓冲\\慢消费者};
\draw[->] (a)--(b);\draw[->] (b)--(c);\draw[->] (c)--(d);
\draw[->,dashed] (d.south) -- ++(0,-.6) -| (a.south);
\node[below=8mm of b,align=center] {容量不足、取消和期限到达向上游传播};
\end{tikzpicture}
\caption[多级请求队列背压反馈回路]{多级请求队列背压反馈回路。各处理层级设定有限容量预算，下游负载与队列释放信号向上游形成背压反馈以防止显存过载。}\label{fig:schedule-backpressure}
\end{figure}

突发到达时，队列吸收的只是暂时差额，永久性不足仍需扩容解决。若一段时间内输入速率为 $\lambda_b$、有效完成速率为 $\mu_b<\lambda_b$，持续 $\tau$ 秒，则新增积压粗估为 $(\lambda_b-\mu_b)\tau$。这是流体近似，忽略随机波动；突发形状和请求工作量另需计入。扩容有加载延迟时，尚未就绪的设备尚未构成即时服务能力。

\section{抢占调度}
\optional
\subsection{公平调度的分配单位}
\term{先到先服务}{First-Come, First-Served}{FCFS}简单且可解释，但长请求可能造成\term{队首阻塞}{Head-of-Line Blocking}{HOL Blocking}。按最短作业优先在已知独立作业时长等条件下改善平均等待，而输出长度未知，估计误差和长作业饥饿会削弱其适用性。

租户公平可以按请求数、词元数或计算成本定义。每请求一票偏向长请求，每词元一票又抹平了长上下文解码与短上下文解码的成本差异。应选可测量、能核算的成本单位，并公开权重的业务含义。

一种实现是为租户 $j$ 维护信用额 $D_j$，每轮增加 $w_jQ$，执行后减去计费成本 $c$。信用不足的候选等待后续轮次；允许有限结转可以容纳较大任务，上限则用来限制积累造成的突发。预测成本用于选择，实际成本用于事后纠偏；二者差额长期由其他租户承担会破坏配额。共享批次成本难以精确归因，可按固定的成本分摊规则核算。

\term{老化}{Aging}{}随等待增加优先级，可缓解饥饿；持续超载时，到达工作量超过服务能力，任何优先级策略都会积累无法按时完成的请求。截止时间优先的经典最优性依赖单处理器等前提，本系统还具有缓存驻留、多资源约束和不可抢占设备片段，这些前提并不成立。

\subsection{抢占成本}
\term{抢占}{Preemption}{}应发生在安全的迭代边界。已完成输出保持不变，未提交候选与执行状态按明确规则处理。若丢弃缓存，恢复需\emph{重算}已提交前缀；若\emph{换出}，则需传输及恢复页映射，并确认模型修订、位置和缓存格式一致。

设重算成本为 $C_r$，换出和换入共需移动 $2m$ 字节，有效传输带宽为 $b$，另有同步成本 $C_s$。粗略比较 $C_r$ 与 $\frac{2m}{b}+C_s$ 有助于选择，而链路竞争与恢复时设备占用还会改变两者的实际大小。频繁抢占会产生反复重算，使观测到的忙碌增加而有效完成减少。应记录重算词元与原始词元的比例，并设置冷却或最小执行份额。

\begin{algorithm}[单轮有界调度]
\AlgoInput{等待队列、活动请求、物理页状态、租户信用和本轮预算。}
\begin{enumerate}
\item 在安全边界处理完成、取消和到期事件，释放引用与预留；已发布终止状态不得重新入队。
\item 增加各租户信用，更新活动解码与等待预填充的期限、年龄与成本估计。
\item 为交互解码与等待预填充分别保留最低服务份额；按权重、期限及年龄选择候选。
\item 对候选确定 $q_i$，同时检查新词元、序列槽、物理页、工作区和轮时预算。无法容纳时考虑其他可行候选，避免无限队首阻塞。
\item 对本轮计划原子预留页与状态版本。若预留失败，撤销并回滚该计划，重新选择候选；仅当所有依赖资源均预留成功时方可转入执行。
\item 执行一轮；按实际进度提交缓存和输出，按实际计费规则扣减信用，保留尚未完成者。
\item 记录排队、批形状、资源、重算和输出阻塞；若不存在可行计划，等待资源事件或执行明确的拒绝、抢占策略，禁止空转。
\end{enumerate}
\end{algorithm}

\section{服务容量评估}

请求级 TPOT 平均值会掩盖单次长停顿，也与把所有 ITL 汇总后的均值不同：后者让长输出拥有更大权重。应分别报告请求级分布与词元级分布，并说明权重。输入长度、输出长度、缓存命中、租户和模型修订都可能形成不同的时延切片。

\term{开放环负载}{Open-Loop Load}{}按独立于完成速度的预设时间表提交请求，可直接显露系统在过载下的排队积压；\term{闭合环负载}{Closed-Loop Load}{}由固定并发用户在收到前序响应后方提交下一请求，其到达率随服务端处理延迟增加而自适应下降。若使用闭合环测试模拟外部固定到达流，由于客户端在延迟剧增时停止发送，拥塞期间本应到达的请求将被平抑或遗漏，这种测量偏差称为\term{协调遗漏}{Coordinated Omission}{}。因此，评估系统承载容量与突发抗压能力须基于\emph{开放环负载}及对应的到达过程。

将满足质量与时延目标的成功请求数除以观察时长，可定义该目标下的\term{有效吞吐}{Goodput}{}。例如在时长为 $\qty{100}{\second}$ 的观察窗口内共提交 $1000$ 个请求，系统拒绝 $100$ 个、执行失败 $50$ 个、成功完成 $850$ 个；其中 $700$ 个同时满足预设的质量与 SLO 时延目标。此时请求提交率为 $\qty{10}{\text{req}/\second}$，成功吞吐为 $\qty{8.5}{\text{req}/\second}$，有效吞吐为 $\qty{7}{\text{req}/\second}$。若仅以成功请求为基准计算目标达标率，结果为 $\frac{700}{850} \approx 82.4\%$；若以全部提交请求为基准，实际达标率为 $70\%$。评估服务承载能力时须显式声明分母口径与统计边界。

调度策略比较应保持模型、生成参数、输入轨迹与输出预算一致；随机生成导致工作量不同，应多次重复或使用可解释的受控长度实验。报告原始到达计划、接纳、失败、超时与取消，分别统计完成时延及各类失败比例。输出质量与长度保持一致，以比较调度收益。



% BEGIN GENERATED INTERVIEWS

% END GENERATED INTERVIEWS
\chapterreferences
\myendchapterdocument
